\section*{NeurIPS Paper Checklist}

\begin{enumerate}

\item {\bf Claims}
    \item[] Question: Do the main claims made in the abstract and introduction accurately reflect the paper's contributions and scope?
    \item[] Answer: \answerYes{}
    \item[] Justification: The abstract claims that typed interaction modules are necessary for compositional generalization (supported by SingleModule vs.\ CausalComp ablation in Tables~1--2, capacity analysis in Table~4, and representation transferability in Table~5). The abstract explicitly states we do not claim state-of-the-art absolute prediction accuracy.

\item {\bf Limitations}
    \item[] Question: Does the paper discuss the limitations of the work performed by the authors?
    \item[] Answer: \answerYes{}
    \item[] Justification: Section~5 discusses four specific limitations: (1) synthetic-only environments, (2) number of interaction types as a hyperparameter, (3) edge supervision requiring ground-truth collision events, and (4) GT object states without end-to-end validation.

\item {\bf Theory assumptions and proofs}
    \item[] Question: For each theoretical result, does the paper provide the full set of assumptions and a complete (and correct) proof?
    \item[] Answer: \answerNA{}
    \item[] Justification: The paper provides theoretical motivation (Section~3.5) based on the Independent Causal Mechanisms principle and causal factorization, but does not claim formal theoretical results. The arguments are informal and clearly stated as such.

    \item {\bf Experimental result reproducibility}
    \item[] Question: Does the paper fully disclose all the information needed to reproduce the main experimental results of the paper to the extent that it affects the main claims and/or conclusions of the paper (regardless of whether the code and data are provided or not)?
    \item[] Answer: \answerYes{}
    \item[] Justification: Appendix~A provides full architecture details (layer sizes, activation functions), training details (optimizer, learning rate, schedule, number of epochs), and environment specifications. All synthetic environments are procedurally generated from described parameters.


\item {\bf Open access to data and code}
    \item[] Question: Does the paper provide open access to the data and code, with sufficient instructions to faithfully reproduce the main experimental results, as described in supplemental material?
    \item[] Answer: \answerYes{}
    \item[] Justification: Code will be released as supplementary material. All datasets are procedurally generated (no external data downloads required). The NRI Springs benchmark follows the standard protocol from Kipf et al.\ (2018).


\item {\bf Experimental setting/details}
    \item[] Question: Does the paper specify all the training and test details (e.g., data splits, hyperparameters, how they were chosen, type of optimizer) necessary to understand the results?
    \item[] Answer: \answerYes{}
    \item[] Justification: Section~4.1 describes environments, compositional splits, and metrics. Section~4.4 reports the hyperparameter sweep. Appendix~A provides full training details including optimizer (AdamW), learning rate ($3 \times 10^{-4}$), weight decay, gradient clipping, and number of epochs (150).

\item {\bf Experiment statistical significance}
    \item[] Question: Does the paper report error bars suitably and correctly defined or other appropriate information about the statistical significance of the experiments?
    \item[] Answer: \answerYes{}
    \item[] Justification: All results in Tables~1--3 report mean $\pm$ standard deviation over 3 random seeds. The seeds control both data generation and model initialization. Error bars capture variability from random data splits and parameter initialization.

\item {\bf Experiments compute resources}
    \item[] Question: For each experiment, does the paper provide sufficient information on the computer resources (type of compute workers, memory, time of execution) needed to reproduce the experiments?
    \item[] Answer: \answerYes{}
    \item[] Justification: Experiments were conducted on a single NVIDIA A100 80GB GPU. Each model trains in approximately 30--90 minutes depending on environment complexity. The full evaluation (3 environments $\times$ 5 methods $\times$ 3 seeds) requires approximately 50 GPU-hours total.

\item {\bf Code of ethics}
    \item[] Question: Does the research conducted in the paper conform, in every respect, with the NeurIPS Code of Ethics \url{https://neurips.cc/public/EthicsGuidelines}?
    \item[] Answer: \answerYes{}
    \item[] Justification: This work develops a method for physics prediction in synthetic environments. It does not involve human subjects, personal data, or dual-use concerns.


\item {\bf Broader impacts}
    \item[] Question: Does the paper discuss both potential positive societal impacts and negative societal impacts of the work performed?
    \item[] Answer: \answerNA{}
    \item[] Justification: This is foundational research on world models for physical dynamics prediction. It does not have direct negative societal impacts. Potential positive impacts include improved robot planning and physical reasoning, but these are distant applications.

\item {\bf Safeguards}
    \item[] Question: Does the paper describe safeguards that have been put in place for responsible release of data or models that have a high risk for misuse (e.g., pre-trained language models, image generators, or scraped datasets)?
    \item[] Answer: \answerNA{}
    \item[] Justification: The released models are small physics prediction networks trained on synthetic data. They pose no risk for misuse.

\item {\bf Licenses for existing assets}
    \item[] Question: Are the creators or original owners of assets (e.g., code, data, models), used in the paper, properly credited and are the license and terms of use explicitly mentioned and properly respected?
    \item[] Answer: \answerYes{}
    \item[] Justification: The NRI Springs benchmark is cited (Kipf et al., 2018). Our implementation follows the standard protocol. All other environments and code are original.

\item {\bf New assets}
    \item[] Question: Are new assets introduced in the paper well documented and is the documentation provided alongside the assets?
    \item[] Answer: \answerYes{}
    \item[] Justification: We introduce two new synthetic environments (SimplePhysics, MultiPhysics) and compositional evaluation splits. These are fully documented in Section~4.1 and Appendix~A, with code provided as supplementary material.

\item {\bf Crowdsourcing and research with human subjects}
    \item[] Question: For crowdsourcing experiments and research with human subjects, does the paper include the full text of instructions given to participants and screenshots, if applicable, as well as details about compensation (if any)?
    \item[] Answer: \answerNA{}
    \item[] Justification: This paper does not involve crowdsourcing or research with human subjects.

\item {\bf Institutional review board (IRB) approvals or equivalent for research with human subjects}
    \item[] Question: Does the paper describe potential risks incurred by study participants, whether such risks were disclosed to the subjects, and whether Institutional Review Board (IRB) approvals (or an equivalent approval/review based on the requirements of your country or institution) were obtained?
    \item[] Answer: \answerNA{}
    \item[] Justification: This paper does not involve research with human subjects.

\item {\bf Declaration of LLM usage}
    \item[] Question: Does the paper describe the usage of LLMs if it is an important, original, or non-standard component of the core methods in this research? Note that if the LLM is used only for writing, editing, or formatting purposes and does \emph{not} impact the core methodology, scientific rigor, or originality of the research, declaration is not required.
    \item[] Answer: \answerNA{}
    \item[] Justification: LLMs are not used as a component of the core methodology. The method operates on object state vectors with MLPs and GNNs, without any language model component.

\end{enumerate}
